%%%%%%%%%%%%%%%%%%%%%%%%%%%%%%%%%%%%%%%%%%%%%%%%%%%%%%%%%%
% # General Discussion and Conclusion
%%%%%%%%%%%%%%%%%%%%%%%%%%%%%%%%%%%%%%%%%%%%%%%%%%%%%%%%%%

% \excnt=122

\chapter{General Discussion and Conclusion}\label{ch-discussion}

Now that we have examined the results from each of the experiments, we are in a position to assess the degree to which the dissertation's objectives have been met. I begin in Section~\ref{dis-summary} by summarizing the results from the experiments and commenting on their broader implications. Section~\ref{dis-implication} then reviews the main findings as they pertain to each of the central theoretical issues explored in the dissertation, and Section~\ref{dis-applications} considers a few additional ways in which this project might be of use to language researchers as well as to the wider community. In Section~\ref{dis-limitation}, I go on to discuss the project's limitations and to make suggestions for future research. Section~\ref{dis-conclusion} concludes the dissertation.

%========================================================%
\section{Summary of the main findings}\label{dis-summary}
%========================================================%

%--------------------------------------------------------%
\subsection{Processing effects for resumption in RC production and comprehension}\label{dis-processing}
%--------------------------------------------------------%

The data presented in this dissertation provide clear-cut evidence that resumption facilitates both production and comprehension of English ORCs and SRCs. In the EPT, rates of resumption responses increased down the SLW hierarchy (from short-distance to long-distance to \textit{wh}-island) for all groups (see~Section~\ref{ept-result-token}). They were most common in the Island condition, where they overtook gap responses for all groups in the SRC dataset. This overall pattern of results, with resumption responses increasing in frequency as the environment grows more difficult for relativization, indicates that resumption eases production of English RCs under processing strain. The fact that the same pattern was observed even with RP acceptors\footnotemark{} removed (see~Section~\ref{rnl-ept-token}) confirms that this was a processing effect rather than a purely representational one. These findings add to the existing literature showing that resumption can assist with RC production both among ENSs (\cite{ferreira.swets2005}; \cite{morgan.wagers2018}; \cite{zukowski.larsen2004}) and among L2ers (\cite{algady2013}; \cite{hyltenstam1984,hyltenstam1990}; \cite{kim2013}; \cite{maghrabi1997}; \cite{pavesi1986}).

\footnotetext{As a reminder, participants were categorized as RP acceptors if they gave `acceptable' ratings (4, 5, or 6) to ≥~80\% of RP trials in at least one condition; conversely, participants were categorized as RP rejectors if they gave `unacceptable' ratings (1, 2, or 3) to ≥~80\% of RP trials in all the RC environments tested.}

In the SPRT, analyses of region-of-interest RTs (Section~\ref{spr-rt}), comprehension-question accuracy scores (Section~\ref{spr-accuracy}), and comprehension-question RTs (Section~\ref{spr-qrt}) all pointed toward resumption making RCs easier to comprehend. The size of the RP-over-gap advantage tended to be larger toward the bottom of the SLW hierarchy, especially in the SRC dataset, indicating that resumption is most helpful where processing load is greatest. Again, the general pattern observed in the data held even with RP acceptors removed (see~Sections~\ref{rnl-spr-rt} and \ref{rnl-spr-accuracy}), indicating that one need not consider resumptive RCs acceptable to derive a processing benefit from them. These results contribute to a growing body of evidence that argues for resumption being able to facilitate RC comprehension (\cite{beltrama.xiang2016}; \cite{hammerly2022}; \cite{hofmeister.norcliffe2013}; \cite{zenker.schwartz2021a}). The SPRT conducted in this dissertation also expands upon previous research by directly testing whether participants accurately interpret resumptive RCs; prior to the present project, only a couple of studies (\cite{hammerly2022}; \cite{morgan.etal2020}) had done this, and neither had found that resumption makes it easier for comprehenders to arrive at the correct interpretation.

At first glance, my SPRT findings might appear to conflict with those from Morgan et al.'s (\citeyear{morgan.etal2020}) study, which found that RPs hinder RC comprehension.\footnotemark{} However, as discussed in Section~\ref{res-morgan}, there is a good reason that Morgan et al.'s data patterned so differently: Their stimuli contained an intervener that was a plausible (feature-matching) antecedent for the pronoun, thereby completely nullifying any processing benefit that might otherwise have been conferred by resumption. In (\ref{ex-piggy1}), for instance, the intervener \textit{Miss Cat} is a perfect fit for the singular and feminine features on the pronoun \textit{her}. 

\footnotetext{Note that each stimulus in the \textcite{morgan.etal2020} study was a cleft sentence, which contains an RC.}

\ex\label{ex-piggy1} 
\ljudge*It was Miss Piggy that Miss Cat reported that Mr.~Dog poked her with a pencil.\par\nobreak\vspace{5pt}
\hfill (\cite{morgan.etal2020}, p. 6, Table 3)
\xe 

Although Morgan et al. may have been incorrect in arguing that RPs do not assist with RC comprehension, their findings are still of great theoretical value because they suggest that RPs behave just like plain pronouns in triggering a search for a viable antecedent (\cite{doron2011}, p.~290; \cite{prince1990}, p.~492; \cite{sells1987}, p.~292). If in a resumptive RC there is another nominal between the head noun and the pronoun that is a potential (feature-matching) antecedent for the pronoun (thereby making it the local antecedent option), resumption can indeed impede dependency resolution instead of facilitating it (see~also \cite{keshev.etal2019}).

Indeed, it would be very strange if RPs did not behave like plain pronouns because they are morphologically and lexically identical to plain pronouns across all documented natural languages (\cite{asudeh2015}, p.~10; \cite{engdahl1982}, p.~172; \cite{hewett2023}, p.~4ff; \cite{mccloskey2017}, p.~5). What is more, there are times when it is impossible for the comprehender to know for sure whether a pronoun in an RC is a plain pronoun or an RP at the point where it is encountered. Take for example (\ref{ex-piggy2}), which is identical to (\ref{ex-piggy1}) up to and including the pronoun \textit{her}.

\ex\label{ex-piggy2}
It was Miss Piggy that Miss Cat reported that Mr.~Dog poked her in front of \_\_.
\xe 

\noindent Unlike in (\ref{ex-piggy1}), having the intervener be the antecedent of the pronoun turns out to be the right choice in (\ref{ex-piggy2}) because there is a gap later in the sentence that must serve as the RC's foot. The pronoun and the gap cannot have the same antecedent because a (nonreflexive) pronoun cannot be coreferential with another argument in the same argument grid (e.g.,~\cite{ogrady2005}, p.~39). Within the sentence, \textit{Miss Cat} is therefore the only possible antecedent for \textit{her}. In (\ref{ex-piggy1}), on the other hand, there is no other potential foot for the RC besides the pronoun, and so we are ultimately forced to interpret it is an RP or else abandon the dependency.

\footnotetext{excluding copular clauses like \textit{He is John} (\cite{haegeman1994}, p.~229; \cite{ogrady2005}, p.~35)}

Modest L2 proficiency effects were observed in both the EPT and the SPRT. In the EPT response data, resumption rates decreased with rising proficiency, at least in the ORC dataset. In the SPRT, RTs in general got faster as proficiency scores went up (in the ORC dataset only), but the effect of proficiency on the size of the [RP~−~gap] contrast did not approach significance. For comprehension-question accuracy in the SPRT, accuracy in general increased as proficiency scores got higher (in both the ORC data and the SRC data), but once again the effect of proficiency on the size of the [RP~−~gap] contrast did not approach significance. These proficiency results indicate both (a)~that as processing efficiency increases, L2ers have less need to resort to resumption when producing English RCs and (b)~that in production as well as in comprehension, higher-proficiency L2ers process English RCs more efficiently than lower-proficiency L2ers.

For ENSs, evidence that resumption facilitates processing during RC production and comprehension was largely restricted to the Island condition, and this was expected since the Long-condition stimuli I used were much less complex than the long-distance conditions in previous adult L1 English studies (e.g.,~\cite{hofmeister.norcliffe2013}). My prediction that the Long condition would place significant strain on the processor only for L2ers (see~Section~\ref{stu-slw}) was therefore borne out in the data analyses. This finding is consistent with previous research indicating that sentence processing in general requires more effort in a nonnative language than in a native language (\cite{hopp2014,kilborn1992}). My use of relatively easy Long-condition stimuli was an important methodological strength of this project because it helped to put the contrast between L1 and L2 processing efficiency on full display.

Taken together, the EPT and SPRT findings unequivocally argue against the selfish speaker model described in Section~\ref{res-selfish} (\cite{ferreira.swets2005,heestand.etal2011,morgan.etal2020,polinsky.etal2013}) by providing strong evidence that resumption in RCs facilitates not only production but also comprehension, at least when (a)~there is strain on the processor and (b)~the resumptive nominal uniquely identifies the head noun. For some advocates of the selfish speaker model, data suggesting that resumption assists with production but not comprehension is taken as a sign that there is a deep rift between language production and language comprehension, with the two systems consulting morphosyntactic representations that are at least partially distinct from one another (e.g.,~\cite{ferreira.swets2005}, p.~277; \cite{morgan.etal2020}, p.~19). However, my findings cast serious doubt on the idea that any such rift exists. 

% By this logic, my findings suggest that the two systems are in fact intimately connected to each other. The present dissertation, in this respect, challenges traditional accounts of language which maintain that there are distinct processing channels for production and comprehension (\cite{benshalom.poeppel2008}, p.~120; \cite{ferreira.swets2005}, p.~276; \cite{levelt.etal1999}, p.~35; see~\cite{pickering.garrod2013}, p.~329) and contributes an additional piece of evidence suggesting that, setting aside articulatory and auditory processes, production and comprehension are tightly interwoven systems that draw on many of the same cognitive resources (\cite{branigan.etal2000}, p.~B22; \cite{cutler.norris1999}, p.~40; \cite{momma.phillips2018}, p.~248; \cite{ogrady2005}, p.~7; \cite{ogrady2022a}, p.~33; \cite{pickering.garrod2013}, p.~346; \cite{sag.wasow1999}, p.~224; \cite{sharwoodsmith.truscott2014}, p.~96).

%--------------------------------------------------------%
\subsection{L1 and L2 Acceptability of resumptive RCs in English}\label{dis-acceptability}
%--------------------------------------------------------%

As a group, ENSs gave resumptive RCs very low ratings in the English AJT (see~Section~\ref{ajt-rating}). Although their RP ratings did rise somewhat down the SLW hierarchy, even going so far as to surpass gap ratings in the SRC Island condition, the fact remains that their mean RP ratings remained securely in the lower half of the rating scale across all conditions. I therefore interpret this slight upward trend as a sign of processing facilitation for illicit RPs relative to illicit gaps rather than as an indication that ENSs considered Island-condition RPs acceptable (see~also \cite{han.etal2012}, p.~158; \cite{mcdaniel.cowart1999}, p.~B16). After all, only a relatively small number of ENSs consistently accepted RPs even in the Island condition (see~Section~\ref{ajt-individual}). These results corroborate other studies which have shown that ENSs tend to give very low acceptability ratings to sentences containing RPs (\cite{alexopoulou.keller2007}; \cite{han.etal2012}; \cite{heestand.etal2011}; \cite{keffala.goodall2011}; \cite{mcdaniel.cowart1999}), thereby indicating that they regard such sentences as ill-formed.

The L2 results in the English AJT were more of a mixed bag. Compared to ENSs, KLEs and MLEs gave higher ratings on average to trials containing resumptive RCs (see~Section~\ref{ajt-rating}). In fact, the mean RP ratings for KLEs and MLEs crossed over into the upper half of the rating scale in at least the Island condition, if not also the Short and Long conditions (MLEs in the SRC dataset only), suggesting that a sizable portion of the L2ers consistently rated RPs as acceptable. Examination of the by-participant rating patterns (see~Section~\ref{ajt-individual}) confirmed that there were large numbers of consistent RP acceptors in the KLE and MLE groups. Rates of RP acceptance tended to be higher among MLEs than KLEs, particularly in the SRC Long and Island conditions, where more than 50\% of MLEs were categorized as RP acceptors.

On the assumption that grammatical resumption in RCs can be transferred from the L1 to the interlanguage (\cite{hyltenstam1984,hyltenstam1990}; see also \cite{schwartz.sprouse1994,schwartz.sprouse1996}), the presence of at least some RP acceptors in the MLE group was expected because Mandarin is thought to allow resumption in all but one of the ORC/SRC environments tested, i.e.,~short-distance SRCs (\cite{keenan.comrie1977}, p.~93; \cite{li.wei2014}, p.~286; \cite{pan2016}, p.~292ff). High rates of RP acceptance were not expected in the KLE group, however, because Korean is thought to disallow resumption across all these environments (\cite{keenan.comrie1977}, p.~93; \cite{han2013}, p.~339). Further analysis of the by-participant rating data showed that many of the KLEs and MLEs who accepted RPs in English also accepted them in the same environments in their L1 (see~Section~\ref{ajt-transfer}), suggesting that they may have transferred grammatical resumption from the L1 to the interlanguage. However, roughly a quarter of the L2ers I tested displayed RP acceptance that could not be straightforwardly attributed either to L1 transfer or to a general \textit{yes} bias, indicating that there is another avenue besides L1 transfer and target-language input by which L2ers can come to allow RPs in their interlanguage English (as originally suggested by \cite{hyltenstam1984}). While the finding that there is this third category of such L2ers who consistently accept resumptive RCs certainly seems to suggest that it represents a licit option for relativization in their interlanguage grammars, further testing could help to shed more light on why they rated sentences the way they did---for discussion, see~Section~\ref{dis-limitation-ajt} below.

L2 proficiency effects were also observed in the English AJT data such that higher English proficiency scores were associated with closer approximation of native-speaker performance. In the analysis of the by-group mean ratings, a clear gap-over-RP preference in RC ratings emerged for both KLEs and MLEs as English proficiency scores got higher (see~Section~\ref{ajt-rating-proficiency}), indicating that L2ers can become sensitive to the English prohibition against resumption in RCs. Similarly, in the analysis of the by-participant response patterns, RP acceptors tended to have lower English proficiency scores than RP rejectors\footnotemark{} did (see~Section~\ref{ajt-individual-proficiency}), providing further evidence that L2ers can come to disallow English resumptive RCs even if they accept them at an earlier point in development. What is more, all L2 proficiency effects in the AJT were (a)~strongest in the short-distance environment, where processing is easiest and English has gapping, and (b)~weak or nonexistent in the \textit{wh}-island environment, where processing is hardest and English bars both gapping and resumption.

\footnotetext{See fn. 2.}

%--------------------------------------------------------%
\subsection{Acceptability of resumptive RCs in L1 Korean and L1 Mandarin}\label{dis-kozh}
%--------------------------------------------------------%

Ratings for RPs in the Korean AJT were somewhat higher than I had expected based on the judgments of my language consultants and most of the existing literature touching on resumption in Korean RCs (e.g.,~\cite{han2013}, p.~339; \cite{keenan.comrie1977}, p.~93; \cite{kim2013}, p.~91), according to which RPs should be disallowed in all the environments tested in this dissertation. Although the KLE group did not consistently accept RPs in any Korean RC environment (see~Section~\ref{ajt-rating}), the 95\% confidence intervals for their RP ratings straddled the midpoint of the rating scale in the ORC `Island' condition, the SRC Long condition, and the SRC `Island' condition, suggesting that individual participants may have differed in how they rated the sentences. Inspection of the by-participant rating patterns in those environments (see~Section~\ref{ajt-individual}) confirmed that there were roughly equal numbers of RP acceptors and RP rejectors.

In the Mandarin AJT, by contrast, ratings for resumptive RCs were somewhat lower than I~had expected based on the judgments of my language consultants and most of the relevant literature (e.g.,~\cite{keenan.comrie1977}, p.~93; \cite{li.wei2014}, p.~286; \cite{pan2016}, p.~292ff), according to which RPs should be permitted in all the RC environments I tested except the SRC Short condition. In the ORC dataset, the 95\% confidence intervals for resumptive RC ratings were on the lower half of the rating scale in all three environments, indicating that MLEs as a group tended to reject RPs there. The by-group RP ratings were closer to my expectations in the SRC dataset, with the 95\% confidence intervals below the midpoint of the scale in the Short condition, where I had anticipated that MLEs would reject Mandarin RPs, but above the midpoint of the scale in the Long and `Island' conditions, where had I anticipated they would accept Mandarin RPs. Inspection of the by-participant rating patterns showed that only about half of the MLEs consistently accepted RPs in the SRC Long and `Island' conditions, however, showing that there was a lack of consensus regarding how to rate RPs even in those contexts.\footnotemark{}

\footnotetext{There were also MLEs who consistently rejected RPs in these same environments, albeit at much lower proportions: 23\% in the SRC Long condition and 22\% in the SRC `Island' condition.}

One key prediction that was borne out was that the KLE group tended to accept gapped RCs in the Korean `Island' condition. In both the ORC data and the SRC data, the 95\% confidence intervals in that condition were above the midpoint on the rating scale and roughly on par with those in the Long condition. This finding confirms my expectation that because Korean lacks filler-gap RCs that extend into an embedded question introduced by a moved \textit{wh}-phrase (see~Sections~\ref{gap-lw} and \ref{acq-slw-pathway}), Korean renditions of the English Island-condition stimuli would be treated the same as Korean renditions of the English Long-condition stimuli. However, this pattern of results did not hold in the Mandarin AJT data, even though I expected that Mandarin would also lack \textit{wh}-island effects in gapped RCs (again, see Sections~\ref{gap-lw} and \ref{acq-slw-pathway}). Instead, the MLE group on average rated gaps much lower in the Mandarin `Island' condition than in the Mandarin Long condition, and the 95\% confidence intervals for their gap ratings in the `Island' condition either landed on the lower half of the scale (in~the ORC dataset) or straddled the midpoint of the scale (in~the SRC dataset). Thus, only KLEs and not MLEs tended to accept `Island'-condition gapped RCs when tested in their L1. 

Follow-up conversations with my Mandarin consultants have illuminated a few potential causes of the unexpectedly low mean ratings for certain types of gapped and resumptive RCs in Mandarin. One possibility for the disparity between my predictions and the observed patterns of results is that in trying to keep the Mandarin (and Korean) stimuli as close as possible to their English counterparts, I may have failed to make them as natural-sounding as they could be. For example, a possible reason that MLEs did not consistently accept resumptive ORCs in the Mandarin AJT could be that the RP was too close to the head noun. As shown in (\ref{ex-dis-mandarin-orc})---an~example item in the ORC `Island' condition---the RP and the head noun were separated only by the adnominal particle \textit{de}.

\ex[everygla=\zh]
\begingl
\glpreamble `Island' environment in the Mandarin ORC acceptability judgment task: // \label{ex-dis-mandarin-orc}
\gla 黃 麗 認識 那 個 我 好奇 哪 些 偵探 上 個 星期 逮捕 \nogloss{\{\_\_/} @ \boxing{他}\} 的 \boxing{男人}。// 
\glb Huang Li renshi na ge wo haoqi na xie zhentan shang ge xingqi daibu ta de nanren //
\glc Huang Li know {\sc dist} {\sc cl} I curious which {\sc cl} detective last {\sc cl} week arrest {\sc 3ms} {\sc adn} man //
\glft `Huang Li knows the man that I wonder which detectives arrested \{*\_\_/*him\} last week.' //
\endgl 
\xe

\noindent One of my Mandarin consultants informed me that this [RP~+~\textit{de}~+~head~noun] configuration is infelicitous because \textit{de} can be a possessive marker in addition to a relativizer, and thus the string \textit{ta de nanren} can be interpreted locally as \textit{his man}. Present across all three ORC environments in RP trials, this [RP~+~\textit{de}~+~head~noun] sequence could help to account for the low overall ratings for resumptive ORCs. The same language consultant also speculated on the reason for the low gap ratings in the Mandarin ORC/SRC `Island' conditions: Co-occurrence of the [determiner~+~classifier] sequence (i.e.,~\textit{na ge}) and the \textit{wh}-phrase (i.e.,~\textit{na xie zhentan}) in the same RC may cause parsing difficulty for some participants; the reason for this is that \textit{na ge} initiates a search for a head noun, and this search is disrupted by the \textit{wh}-phrase, which renders the clause an embedded question and thus introduces more complexity into the sentence than an ordinary intervening NP would. If this idea is on the right track, then simply omitting the optional \textit{na ge} at the onset of the RC might lead to higher ratings in those conditions.\footnotemark{} It has also been suggested to me (p.c.~Li `Julie' Jiang, 10 May 2024) that ratings in the Mandarin ORC/SRC `Island' conditions might have been higher if I had translated wonder as \textit{xiang zhidao} (`want to know') instead of \textit{haoqi} (`curious')---cf.~Chapter~\ref{ch-ajt}, fn.~1.

\footnotetext{My reason for including \textit{na ge} in the first place was to signal the onset of the RC (recall that Mandarin has head-final RCs), which I had hoped would cut down on garden-pathing (see~\cite{lin.bever2006}, p.~255), and make the sentence as a whole easier to process.}

To summarize, the Korean and Mandarin AJT results deviated from my expectations and the extant literature in several ways. First, RP ratings were higher in Korean and lower in Mandarin than I had predicted. Second, gap ratings in the Mandarin `Island' condition were lower than I had expected. Finally, rating patterns were far less uniform in L1 Korean and L1 Mandarin than in L1 English. These results highlight the need for more experimental work on gapped and resumptive RCs in both Korean and Mandarin.

%--------------------------------------------------------%
\subsection{Relationship between acceptability and processing effects}\label{dis-assumption}
%--------------------------------------------------------%

In Chapter~\ref{ch-reanalysis}, I investigated whether it was the consistent RP acceptors who were chiefly responsible for the apparent facilitation effects for resumptive RCs found in the group-level analyses of the processing data. To do this, I repeated the main EPT and SPRT analyses with RP acceptors in the English AJT removed. Despite a considerable reduction in statistical power, there were still clear signs of processing facilitation in the EPT production data (see~Section~\ref{rnl-ept-token}), the SPRT region-of-interest RT data (see~Section~\ref{rnl-spr-rt}), and the SPRT comprehension-question accuracy data (see~Section~\ref{rnl-spr-accuracy}). 

These results directly refute the assumption traditionally held among L2 researchers that L2ers' production patterns reflect representations licensed in their interlanguage grammars (see~Section~\ref{acq-assumption}). This assumption, which I call the Performance-Competence Correspondence Assumption (PCCA), has often prevailed through the decades in studies looking at L2 gapped vs. resumptive RCs (\cite{algady2013}; \cite{gass1979}; \cite{hyltenstam1984}; \cite{marefat.abdollahnejad2014}; \cite{pavesi1986}). It is my hope that the present findings will help to put this assumption to rest.

Beyond challenging the PCCA, the fact that there were processing effects for resumption even with RP acceptors removed indicates that for at least some of the L2ers I tested, resumptive RCs in their interlanguage English can best be characterized as processing resumptives (\cite{asudeh2004}, p.~295) rather than grammatical resumptives. That is to say, they regard resumptive RCs as ill-formed and yet derive a processing benefit from them in certain domains that are difficult for relativization, just as the present dissertation and previous adult L1 research have found to be the case for ENSs (\cite{chacon2015}, p.~89; \cite{kroch1981}, p.~131ff; \cite{sells1984}, p.~461).

%========================================================%
\section{Central theoretical issues revisited}\label{dis-implication}
%========================================================%

Before moving on to explore the practical applications of the present findings, let us briefly review the main theoretical implications of this research:

%--------------------------------------------------------%
\subsection{Selfish speaker model}
%--------------------------------------------------------%

The EPT and SPRT results indicate that resumption facilitates both production and comprehension of RCs when (a)~there is strain on the processor and (b)~the resumptive nominal uniquely identifies the head noun. These findings directly challenge the selfish speaker model (\cite{ferreira.swets2005}; \cite{heestand.etal2011}; \cite{morgan.etal2020}; \cite{polinsky.etal2013}); in doing so, they challenge the hypothesis that there is only partial overlap among the morphosyntactic representations consulted during language production vs. language comprehension.

%--------------------------------------------------------%
\subsection{Status of resumptive RCs in English interlanguages}
%--------------------------------------------------------%

Some L2ers consistently accepted RPs, indicating that their interlanguage grammars may allow resumption in certain RC environments, whereas other L2ers consistently rejected RPs, indicating that they do not consider resumption a licit option for English relativization. Roughly a quarter of the L2 participants displayed consistent RP acceptance in English that could not be directly attributed either to L1 transfer or to a general \textit{yes} bias, suggesting that there is another way besides L1 transfer and target-language input that can lead L2ers to come to allow RPs in interlanguage English. The source of this RP acceptance remains unclear; still, the findings are \textit{prima facie} evidence for Hyltenstam's (1984, 1990) hypothesis that resumptive RCs enter interlanguage grammars by dint of being relatively easy to process. It will be important to follow up on the present research by conducting additional studies aimed at investigating why some L2ers accept resumptive RCs in such environments (see~Section \ref{dis-limitation-ajt} below).

%--------------------------------------------------------%
\subsection{PCCA}
%--------------------------------------------------------%

Signs of L2 processing facilitation for resumption were detected in the EPT and SPRT data even after RP acceptors were removed. These findings constitute evidence against the traditional L2 assumption, which remains widespread in the L2 literature, that production patterns faithfully reflect representations licensed by the interlanguage grammar (\cite{algady2013}; \cite{gass1979}; \cite{hyltenstam1984}; \cite{marefat.abdollahnejad2014}; \cite{pavesi1986}).

%--------------------------------------------------------%
\subsection{Purging resumption from the interlanguage}
%--------------------------------------------------------%

A proficiency effect was observed in the English AJT data whereby higher L2 proficiency scores were associated with both higher gap ratings and lower RP ratings, at least in environments where gapped RCs are permitted in English. In addition, analysis of the by-participant rating data showed that RP rejectors had higher L2 proficiency scores on average than RP acceptors did. These results indicate that KLEs and MLEs are indeed able to develop a sensitivity to English's prohibition on resumption in RCs, which is consistent with the `overwriting' approach to L2 acquisition of RCs presented in Section~\ref{acq-initial-state}. 

%========================================================%
\section{Practical applications}\label{dis-applications}
%========================================================%

%--------------------------------------------------------%
\subsection{Direct applications}
%--------------------------------------------------------%

There are several other ways in which I believe the information presented in this dissertation could benefit language researchers and the wider community. First and foremost, this dissertation draws on innovative methods honed over the past two decades in the adult L1 English processing literature to gain a deeper understanding of the phenomenon of RC resumption in adult L1/L2 English, and I hope that this project inspires other L2 researchers to leverage psycholinguistic methodologies to explore issues at the intersection of L2 competence and L2 performance.

Second, the present findings could potentially be of use to the English language teacher because they provide insights into why resumptive RCs crop up from time to time in L2ers' spoken and written output. Specifically, my data analyses showed that L2 resumption is not necessarily indicative of nontarget representations. Instead, resumptive RCs are often produced as a means of managing processing load where it is greatest, even among L2ers who do not consider them acceptable in English. Furthermore, the results indicate that L2ers gradually begin to reject RPs as their level of English proficiency rises, at least in RC environments where gapping is relatively easy on the processor.

Finally, I hope that the C-test I created and validated as part of this project will be of use to EFL/ESL programs, where quick assessment of English proficiency is often used for placement purposes. C-tests offer a particularly fast and easy method of assessing a nonnative speaker's general proficiency level in the target language (\cite{bachman1985}; \cite{dornyei.katona1992}; \cite{eckes.grotjahn2006}; \cite{park.choi2018}). Furthermore, I anticipate that the R code I developed for automatic scoring of the C-test responses (see Appendix~\ref{app-validation}, Section~\ref{app-validation-scoring} for details) will help others save time when scoring not only my C-test but also their own C-tests and cloze tests.

%--------------------------------------------------------%
\subsection{Indirect applications}
%--------------------------------------------------------%

There are other, less immediate impacts that research of the kind reported in this dissertation can have on language teaching and society at large. Many advances in technology and medicine grow out of theoretical research that had no obvious applications at the time it was carried out. Quantum mechanics was initially seen as a purely theoretical enterprise, and yet it has led to the development of the semiconductors used in modern computers and smartphones (\cite{lenz2016}). By the same token, even the most abstract of linguistic investigations can eventually lead to innovations in pedagogical practice, natural language processing, artificial intelligence, and the like. For instance, many of the teaching techniques that are now common in L2 classrooms have their roots in research that was once considered esoteric (\cite{sadeghi2022}, p.~136).

%========================================================%
\section{Limitations and future directions}\label{dis-limitation}
%========================================================%

%--------------------------------------------------------%
\subsection{Imposing time limits on elicited-production responses}\label{dis-limitation-ept}
%--------------------------------------------------------%

There were several potential methodological problems with the experiments presented in this dissertation besides the issues with the Mandarin AJT stimuli already noted in Section~\ref{dis-kozh}. For instance, my decision to impose a time limit on responses in the EPT may have been counterproductive. As a reminder, participants were permitted only six seconds to give their full response once they pushed the button to start recording.\footnotemark{} I had hoped that forcing participants to answer quickly would (a)~keep the influence of metalinguistic knowledge to a minimum and (b)~increase load on the processor, thereby maximizing the chances of their resorting to resumption when formulating a response. In the \textcite{ferreira.swets2005} RC elicitation task, which I detail in Section~\ref{res-evidence-production}, ENSs produced resumptive RCs at high rates even without a time limit, however, so it is doubtful that my decision to give participants less time to answer had the desired effect of maximizing resumption responses. 

\footnotetext{Participants did have the option to re-record their responses as many times as they wished. Including this `re-record' option was likely also a misstep, however, because it meant that participants may have been overwriting their initial (and potentially most informative) responses at least some of the time.}

Furthermore, there are two potential problems associated with time limits that researchers should keep in mind when designing EPTs like the one reported in this dissertation. First, some recordings were cut off before participants finished giving their response, leading to loss of a small amount of the EPT data. Such responses were coded as Nontarget and excluded from further analysis. Second, forcing participants to squeeze their responses into a narrow time window made an RT analysis of their spoken responses effectively pointless. \textcite{ferreira.swets2005}, who imposed no such limit on responses once participants had initiated their recordings, found that toward the end of the sentence, RTs in their critical condition eliciting \textit{wh}-island SRCs, as in (\ref{ex-dis-ferreira-a}), were significantly slower than RTs in their control condition eliciting short-distance SRCs, as in (\ref{ex-dis-ferreira-b}). It is possible that I, too, might have found a significant gapped RC vs. resumptive RC contrast in RTs had I given participants more time to complete their responses. Such data could have shed further light on how gapping and resumption compare in real-time RC production.

\pex
\a \ljudge*This is a donkey that I don't know where \boxing{it} lives.\label{ex-dis-ferreira-a}
\a This is a donkey that \boxing{\_\_} doesn't know where it lives.\label{ex-dis-ferreira-b}
\par\nobreak\vspace{5pt}
\hfill (adapted from \cite{ferreira.swets2005}, p. 271, (3))
\xe 

%--------------------------------------------------------%
\subsection{Exploring other methodologies for testing RC comprehension}\label{dis-limitation-spr}
%--------------------------------------------------------%

The present SPRT results strongly indicate that resumption facilitates RC comprehension, at least under certain conditions. However, these findings could arguably be made even more robust through the use of additional experimental methodologies. For instance, Morgan et al.'s (\citeyear{morgan.etal2020}) study probing the comprehension of gapped vs. resumptive RCs included not only an SPRT but also a sentence-picture matching task and a visual-world eye-tracking task. In the sentence-picture matching task, participants needed to click on the picture that best matched the stimulus sentence (Morgan et al., \citeyear{morgan.etal2020}, p.~4). Figure~\ref{fig:dis-spm} shows a sentence-picture matching trial similar to the ones that Morgan et al. used, but with the stimulus sentence (a \textit{whether}-island) altered so that the RP uniquely identifies the head noun.

% SPM Figure

\begin{figure}
\caption{Modified \textit{wh}-island trial from Morgan et al.'s (\citeyear{morgan.etal2020}) sentence-picture matching task} \label{fig:dis-spm}
\raggedright\setstretch{1}
\tikz{
\node[rectangle, draw, minimum width=6.5in, minimum height=3.5cm] (box) at (0,0) {};

\node [font=\sffamily] (stimulus) at ([yshift=-.7cm]box.north) {It is Mr.~Dino that Miss Rabbit wondered whether Miss Piggy tickled him with a feather.};

\node [draw] () at ([yshift=-1.15cm,xshift=-6cm]stimulus.south){\includegraphics[width=3.25cm]{../image/morgan1.png}};

\node [draw] () at ([yshift=-1.15cm,xshift=-2cm]stimulus.south){\includegraphics[width=3.25cm]{../image/morgan2.png}};

\node [draw] () at ([yshift=-1.15cm,xshift=2cm]stimulus.south){\includegraphics[width=3.25cm]{../image/morgan3.png}};

\node [draw] () at ([yshift=-1.15cm,xshift=6cm]stimulus.south){\includegraphics[width=3.25cm]{../image/morgan4.png}};

}
\textit{Note}. Original sentence from Morgan et al.'s study: \textit{It is Mr.~Dino that Mr.~Rabbit wondered whether Miss Piggy tickled him with a feather}; the target picture is the second from the left.
\vspace{-12pt}
\end{figure}

In Morgan et al.'s visual-world eye-tracking task, participants' eye-movements were monitored while they looked at a computer screen and listened to recorded sentences. The visual stimuli contained pictures of four animal characters, one in each corner of the screen. After the recording finished playing, a comprehension question appeared asking participants to identify who did what to whom in the sentence. Figure~\ref{fig:dis-vyp} shows an eye-tracking trial of the kind that Morgan et al. employed in their study, but again with the stimulus sentence (an adjunct \textit{why}-island) modified such that the head noun is the pronoun's only possible antecedent.

% VYP Figure

\begin{figure}
\caption{Modified \textit{wh}-island trial from Morgan et al.'s (\citeyear{morgan.etal2020}) eye-tracking task} \label{fig:dis-vyp}
\raggedright\setstretch{1}
\tikz{
\node[rectangle, draw, minimum width=6.5in, minimum height=6.2cm] (box) at (0,0) {};

\node [font=\sffamily] (stimulus) at ([yshift=-.7cm]box.north) {It was Miss Piggy who Mr.~Cat understood why Mr.~Dog poked her with a pencil.};

\node () at ([xshift=-5pt]stimulus.west){\includegraphics[width=.4cm]{../image/speaker.png}};

\node[rectangle, draw, minimum width=2in, minimum height=1.8in] (world) at ([yshift=-2.5cm]stimulus.south) {};

\node () at ([yshift=1.5cm,xshift=1.6cm]world.center){\includegraphics[height=1.5cm]{../image/dog.png}};

\node () at ([yshift=-1.35cm,xshift=-1.6cm]world.center){\includegraphics[height=1.35cm]{../image/cat.png}};

\node () at ([yshift=-1.35cm,xshift=1.6cm]world.center){\includegraphics[height=1.29cm]{../image/pig.png}};

\node () at ([yshift=1.38cm,xshift=-1.6cm]world.center){\includegraphics[height=1.2cm]{../image/rabbit.png}};

}
\textit{Note}. Original sentence from Morgan et al.'s study: \textit{It was Miss Piggy who Miss Cat understood why Mr.~Dog poked her with a pencil}; the target picture is in the bottom right-hand corner.
\vspace{-12pt}
\end{figure}

\noindent The corresponding comprehension question is shown in Figure~\ref{fig:dis-question}. Like the comprehension questions in my SPRT, arriving at the correct answer depended on accurate resolution of the RC dependency.

% VYP Question

\begin{figure}
\caption{Modified comprehension question from Morgan et al.'s (\citeyear{morgan.etal2020}) eye-tracking task} \label{fig:dis-question}
\raggedright\setstretch{1}
\tikz{
\node[rectangle, draw, minimum width=6.5in, minimum height=3.25cm] (box) at (0,0) {};

\node [font=\sffamily] (stimulus) at ([yshift=-.6cm]box.north) {Who did what do whom?};

\node[rectangle, rounded corners=3pt, draw, line width=1pt, fill=black!10, font=\sffamily\footnotesize, minimum width=7cm, minimum height=.75cm, text width=6.4cm, align=left] (one) at ([yshift=-.6cm,xshift=3.6cm]stimulus.south) {Mr.~Dog poked Miss Piggy with a pencil.};

\node[rectangle, rounded corners=3pt, draw, line width=1pt, fill=black!10, font=\sffamily\footnotesize, minimum width=7cm, minimum height=.75cm, text width=6.4cm, align=left] (two) at ([yshift=-.6cm,xshift=-3.6cm]stimulus.south) {Mr.~Dog poked Mr.~Cat with a pencil.};

\node[rectangle, rounded corners=3pt, draw, line width=1pt, fill=black!10, font=\sffamily\footnotesize, minimum width=7cm, minimum height=.75cm, text width=6.4cm, align=left] (three) at ([yshift=-1.6cm,xshift=-3.6cm]stimulus.south) {Mr.~Dog poked Miss Rabbit with a pencil.};

\node[rectangle, rounded corners=3pt, draw, line width=1pt, fill=black!10, font=\sffamily\footnotesize, minimum width=7cm, minimum height=.75cm, text width=6.4cm, align=left] (four) at ([yshift=-1.6cm,xshift=3.6cm]stimulus.south) {Mr.~Cat poked Miss Piggy with a pencil.};

}
\textit{Note}. The target answer is in the top right-hand corner.
\vspace{-12pt}
\end{figure}

In the response data from their sentence-picture matching task and the comprehension-question accuracy data from their eye-tracking task, Morgan et al. found that there were fewer target responses in RP trials than in gap trials. The effect was strongest in their nonisland condition (where processing should be easiest), as in (\ref{ex-dis-morgan-a}), and weakest in their strong island condition (where processing should be hardest), as in (\ref{ex-dis-morgan-c}), but it was observed at least numerically across all environments tested. Analysis of the eye-movement data revealed that gaps led to increased looks to the head noun (i.e.,~Mr. Dino) whereas RPs resulted in increased looks to both the head noun and the intervener (i.e.,~Mr. Rabbit)---the RP's local antecedent option. 

\pex Conditions from the sentence-picture matching task in \textcite{morgan.etal2020}\par\nobreak
\a Nonisland:\par\nobreak\vspace{5pt}
It is Mr. Dino that Mr. Rabbit said that Miss Piggy tickled \{\_\_/*him\} with a feather.\label{ex-dis-morgan-a}
\a Weak island:\par\nobreak\vspace{5pt}
It is Mr. Dino that Mr. Rabbit wondered whether Miss Piggy tickled \{*\_\_/*him\} with a feather.\label{ex-dis-morgan-b}
\a Strong island:\par\nobreak\vspace{5pt}
It is Mr. Dino that Mr. Rabbit slept while Miss Piggy tickled \{*\_\_/*him\} with a feather.\label{ex-dis-morgan-c}
\par\nobreak\vspace{5pt}
\hfill (adapted from \cite{morgan.etal2020}, p. 5, Table 1)
\xe 

Morgan et al. interpret these results as evidence that RPs interfere with RC comprehension. I strongly suspect, however, that the results would have patterned differently if they had used stimuli like the ones in Figures~\ref{fig:dis-spm} and \ref{fig:dis-vyp}, where the RP is better equipped to assist with dependency resolution. Replicating the present SPRT's findings in the sentence-picture matching and visual-world eye-tracking paradigms (and still other methodologies, such as ERPs, fMRIs, and elicited imitation) could furnish corroborating evidence that resumption is capable of facilitating RC comprehension.

%--------------------------------------------------------%
\subsection{Tapping into well-formedness perceptions}\label{dis-limitation-ajt}
%--------------------------------------------------------%

There is no perfect way to assess perceptions of well-formedness, but AJTs remain widely used in L1 and L2 research because they are the best tool that we have for tapping into those perceptions (e.g.,~\cite{sprouse2018}). In this dissertation, I interpreted consistently (≥~80\%) high AJT ratings as a sign that a sentence type is considered well-formed and consistently (≥~80\%) low AJT ratings as a sign that a sentence type is considered ill-formed. The finding that a sizable portion of my L2ers consistently accepted resumptive RCs in environments where they did not consistently accept them in their L1 therefore suggests that (a)~such L2ers have interlanguage English grammars that allow resumption in those environments and (b)~the source of that grammatically-licensed resumption is something other than L1 transfer. However, there are a number of methodological tweaks one could make that could potentially shed more light on why participants rate sentences the way they do. One possibility would be to have participants correct sentences that they rate as unacceptable (\cite{devilliers.devilliers1972}, p.~303; \cite{hawkins.chan1997}, p.~204; \cite{maghrabi1997}, p.~113). As explained in Section~\ref{acq-hawkins}, \textcite{hawkins.chan1997} found that most of the consistent RP acceptors among their elementary-level L1-Cantonese L2ers of English (CLEs) tended to correct island-gap RCs by inserting an RP. These results were taken as evidence that CLEs initially transfer grammatical resumption from the L1 to the English interlanguage. However, having participants perform an extra task for only sentences that they find unacceptable gives them an incentive to mark sentences as acceptable, and this could have the unintended effect of inflating their ratings.

Think-aloud protocols, whereby participants comment on their thought process as they rate sentences, can offer rich information about participant's judgments (\cite{davies.kaplan1998}; \cite{ericsson.simon1980}; \cite{maghrabi1997}, p.~112ff). \textcite{maghrabi1997} made use of one such think-aloud task in his dissertation examining resumptive RCs among L1-Arabic L2ers of English (ALEs) and L1-English L2ers of Arabic (ELAs). Participants were asked to rate sentences as `grammatical' or `ungrammatical,' correct any errors they found, and explain their choices as they worked through the task. Maghrabi's analysis of the results showed that the main reason ALEs and ELAs gave for accepting resumptive RCs and for adding RPs to gapped RCs was ``to~ease processing load contained within the relative clause constructions in order to understand more easily'' (p.~144). Examples of participants' think-aloud responses are provided in (\ref{ex-dis-maghrabi}).

\newpage

\pex\label{ex-dis-maghrabi}
\a Think-aloud response from ALE \#15 (for a rejected sentence with a gapped RC): \par\nobreak\vspace{5pt}
\dots~should add something to make the sentence clear \dots ~I want to understand what it mean \dots ~what is going on \dots ~you know \dots ~this way no idea I don't understand \dots ~I mean no complete idea \dots
\a Think-aloud response from ELA \#5 (for a rejected sentence with a gapped RC): \par\nobreak\vspace{5pt}
Twenty-six is incorrect \dots{} once again because something is missing \dots{} We need to know we need to know to to who? \dots{} gave my books to the teacher to somebody else? \dots{} because once again the the reader or someone hearing the sentence is left wanting more information \dots{} It's incomplete thought \dots{} left dangling. \par\nobreak\vspace{5pt}
\hfill (\cite{maghrabi1997}, p. 147, Table 18)
\xe 

\noindent Importantly, Maghrabi argues that his think-aloud data show that L2ers who correct gapped RCs by adding an RP may do so not because they regard the RC with the RP as well-formed but instead because they think that it makes the sentence easier to understand. If this were so, however, it might suggest that the participants were not well trained in what they were supposed to do. Of course, a tendency for participants to comment on the relative processing ease of resumptive RCs does not necessarily preclude the possibility that they also consider such RCs to be well-formed.

It may be that some of my L2 participants were more influenced by processing ease than by well-formedness perceptions when they rated resumptive RCs as acceptable. One way to address this potential problem would be to look for ways to modify the instructions to make it clearer what they are expected to do in the task (see~\cite{bley-vroman.etal1988}). Another way would be to add a debriefing session at the end of the experiment where participants are asked pointed questions about why they rated certain trials the way they did. 

It has been suggested that certain types of participants, such as heritage speakers, overaccept illicit items in AJTs out of a ``lack of confidence in their linguistic intuition'' (\cite{rinke.flores2014}, p.~688--689; see also, e.g., \cite{polinsky2016,sherkina-lieber.etal2011}). To assess whether such feelings of insecurity might have inflated L2ers' ratings of, e.g., resumptive RCs, I could to have had them also rate how confident they were in their acceptability judgments. \textcite{dillon.etal2019}, for instance, had ENSs give each sentence a binary grammaticality judgment accompanied by a three-point confidence rating with options ranging from `not at all confident' to `very confident.' A version of this scale adapted to the present study is shown in Figure~\ref{fig:dis-confidence}. Confidence ratings of this type could help to assess how clear participants' well-formedness perceptions are.

\begin{figure}
\caption{Rating scale incorporating both acceptability judgments and confidence judgments} \label{fig:dis-confidence}
\tikz{
\node[rectangle, draw, minimum width=6.5in, minimum height=4cm] (box) at (0,0) {};
\node [font=\sffamily] (stimulus) at ([yshift=-.7cm]box.north) {Mary knows the man that I wonder which book he borrowed last month.};
% \node (accept) at ([yshift=-.5cm,xshift=.5cm]stimulus.south){\includegraphics[width=.7cm]{../image/check.png}};
% \node (reject) at ([yshift=-.5cm,xshift=-.5cm]stimulus.south){\includegraphics[width=.7cm]{../image/remove.png}};
\node (sad2) at ([yshift=-.5cm,xshift=-1.5cm]stimulus.south){\includegraphics[width=.8cm]{../image/sad2.png}};
\node (sad1) at ([yshift=-.5cm,xshift=-.5cm]stimulus.south){\includegraphics[width=.8cm]{../image/sad1.png}};
\node (happy1) at ([yshift=-.5cm,xshift=.5cm]stimulus.south){\includegraphics[width=.8cm]{../image/happy1.png}};
\node (happy2) at ([yshift=-.5cm,xshift=1.5cm]stimulus.south){\includegraphics[width=.8cm]{../image/happy2.png}};
\node[rectangle, rounded corners=3pt, draw, line width=1pt, fill=black!10, font=\sffamily\footnotesize, minimum width=.65cm, minimum height=.65cm] (one) at ([yshift=-1.5cm,xshift=-1.5cm]stimulus.south) {1};
\node[rectangle, rounded corners=3pt, draw, line width=1pt, fill=black!10, font=\sffamily\footnotesize, minimum width=.65cm, minimum height=.65cm] (two) at ([yshift=-1.5cm,xshift=-.5cm]stimulus.south) {2};
\node[rectangle, rounded corners=3pt, draw, line width=1pt, fill=black!10, font=\sffamily\footnotesize, minimum width=.65cm, minimum height=.65cm] (three) at ([yshift=-1.5cm,xshift=.5cm]stimulus.south) {3};
\node[rectangle, rounded corners=3pt, draw, line width=1pt, fill=black!10, font=\sffamily\footnotesize, minimum width=.65cm, minimum height=.65cm] (four) at ([yshift=-1.5cm,xshift=1.5cm]stimulus.south) {4};
\node (least_acceptable) [font=\sffamily\footnotesize, left=0pt of sad2.west, anchor=east] {Completely Unacceptable};
\node (most_acceptable) [font=\sffamily\footnotesize, right=0pt of happy2.east, anchor=west] {Completely Acceptable};
\node (least_confident) [font=\sffamily\footnotesize, left=5pt of one.west, anchor=east] {Least Confident};
\node (most_confident) [font=\sffamily\footnotesize, right=5pt of four.east, anchor=west] {Most Confident};
\node [font=\sffamily\footnotesize] (skip) at ([yshift=-2.35cm]stimulus.south) {Press zero if you cannot rate the sentence.};
}
\vspace{-36pt}
\end{figure}

% Another possibility would be to make the I-don't-know response one of the primary response options. In conventional AJTs, I-don't-know responses are treated as lost data. While it is important not to place the I-don't-know option on the main rating scale (to ensure that participants think of it as separate from the acceptability spectrum), placing it apart from the main scale is also undesirable if we intend to treat I-don't-know responses as valuable data. The reason for this is that the very act of placing the I-don't-know option apart from the main rating scale may mark it as a last-resort option and therefore discourage participants from using it. One way to get around this potential problem would be to offer participants three response options---`acceptable,' `unacceptable,' and `not sure'---that are arranged in a non-linear pattern such as a triangle. Figure~\ref{fig:dis-unsure} shows what such a configuration could look like. The data analyses could then examine ratings along two distinct dimensions: an `acceptability' dimension examining the proportion of trials judged to be acceptable vs. unacceptable and a `ratability' dimension examining the proportion of trials rated vs. skipped.

% \begin{figure}
% \caption{Rating scale for three-way acceptable vs. unacceptable vs. not sure judgments} \label{fig:dis-unsure}
% \tikz{
% \node[rectangle, draw, minimum width=6.5in, minimum height=3.2cm] (box) at (0,0) {};
% \node [font=\sffamily] (stimulus) at ([yshift=-.7cm]box.north) {Mary knows the man that I wonder which book he borrowed last month.};
% \node (accept) at ([yshift=-.6cm]stimulus.south){\includegraphics[width=.7cm]{../image/check.png}};
% \node (reject) at ([yshift=-.8cm,xshift=-.6cm]accept.center){\includegraphics[width=.7cm]{../image/remove.png}};
% \node (unsure) at ([yshift=-.8cm,xshift=.6cm]accept.center){\includegraphics[width=.7cm]{../image/question.png}};
% }
% \vspace{-30pt}
% \end{figure}

% \noindent Instituting one of these alternative rating procedures in the present study might have made it easier to determine whether or not L2ers have clear perceptions of well-formedness for sentences in a given condition.

%--------------------------------------------------------%
\subsection{Examining different pairings of native language and target language}\label{dis-limitation-distance}
%--------------------------------------------------------%

Future studies could also investigate how crosslinguistic differences affect RC processing. As discussed in Section~\ref{pro-compare}, RCs in English, Korean, and Mandarin differ in several notable ways. For instance, English gapped RCs are of the filler-gap type, with the filler identified before the gap, while Korean and Mandarin gapped RCs are of the gap-filler type, with the gap coming before the filler. To understand the way in which differences such as these affect how RCs are processed in real time, let us consider the English short-distance ORC \textit{the book that Lisa bought} and its translation equivalents in Korean and Mandarin, which are shown in (\ref{ex-dis-kor}) and (\ref{ex-dis-zho}), respectively.

\ex[everygla=\ko]
\begingl 
\gla 리사가 산 책 // \label{ex-dis-kor}
\glb Lisa-ka sa-n chayk //
\glc Lisa-{\sc nom} buy-{\sc adn} book //
\glft `the book that Lisa bought' //
\endgl 
\xe 

\ex[everygla=\zh]
\begingl 
\gla 麗莎 買 的 書 // \label{ex-dis-zho}
\glb Lisha mai de shu //
\glc Lisa buy {\sc adn} book //
\glft `the book that Lisa bought' //
\endgl 
\xe

\noindent Processing routines for comprehending these three RCs are provided in Tables \ref{tblr:dis-routine-eng}--\ref{tblr:dis-routine-zho} (using English lexemes for ease of exposition), based on the approach to modeling the real-time processing of RCs presented in Section~\ref{pro-routine}. Notice that while the three routines share many of the same operations, the sequencing of those operations varies across languages, with Korean and Mandarin being much closer to each other than either of them is to English. 

% English

\setstretch{1}
\begin{longtblr}[caption={Processing routine for comprehending sentence-initial \textit{The book that Lisa bought\dots} in English}, label={tblr:dis-routine-eng}, remark{Note}={PRED~=~predicate position; $\alpha$~=~first-argument position; $\beta$~=~second-argument position; vertical lines indicate modification; superscripts are fillers.}]{colspec={|X[c]|}, rowhead = 0}\hline

\begin{tblr}{colspec={X[.6in]X[l]}} 
\textbf{Step 0:} & \textbf{Semantic Base Projection} \vspace{2pt} \\ 
& Project a minimal predicate-argument complex.  \vspace{2pt} \\
& PRED($\alpha$) \\
\end{tblr} \\*

\begin{tikzpicture}
\node[draw=black, minimum height=20pt, minimum width=6.2in] (box) {
{\itshape{\bfseries The book}\dots\phantom{that Lisa bought }}};
\end{tikzpicture} \\*

\begin{tblr}{colspec={X[.6in]X[l]}} 
\textbf{Step 1:} & \textbf{First-Argument Algorithm} \vspace{2pt} \\ 
& Integrate the content of \textit{the book}, represented here as $b$, at the first-argument position.  \vspace{2pt} \\
& PRED($\mathbfit{b}$) \\
\end{tblr} \\

\begin{tikzpicture}
\node[draw=black, minimum height=20pt, minimum width=6.2in] (box) {
{\itshape The book {\bfseries that}\dots\phantom{Lisa bought }}};
\end{tikzpicture} \\*

\begin{tblr}{colspec={X[.6in]X[l]}} 
\textbf{Step 2:} & \textbf{Relativization Algorithm for Head-Initial RCs} \vspace{2pt} \\ 
& Respond to the relativizer \textit{that} by projecting a semantic base as a modifier of $b$ and give it a filler, represented here as a superscript $b$, that needs to be integrated into its argument grid. \vspace{2pt} \\
& 
\begin{tikzpicture}[node distance=0pt, inner sep=0pt]
    \node (1) {\strut PRED(};
    \node (2) [right=of 1.east, anchor=west] {\strut $b$};
    \node (3) [right=of 2.east, anchor=west] {\strut )};
    \node (4) [below=4pt of 2.south, anchor=north] {\bfseries\strut PRED};
    \draw[semithick] ([yshift=2pt]2.south) -- (4.north);
    \node (5) [right=of 4.east, anchor=west] {\bfseries\strut $^{\mathbfit{b}}$(${\mathbfit{\alpha}}$)};
\end{tikzpicture} 
\\
\end{tblr} \\

\begin{tikzpicture}
\node[draw=black, minimum height=20pt, minimum width=6.2in] (box) {
{\itshape The book that {\bfseries Lisa}\dots\phantom{bought }}};
\end{tikzpicture} \\*

\begin{tblr}{colspec={X[.6in]X[l]}} 
\textbf{Step 3:} & \textbf{First-Argument Algorithm} \vspace{2pt} \\ 
& Integrate the content of \textit{Lisa}, represented here as $l$, at the first-argument position. \vspace{2pt} \\
& 
\begin{tikzpicture}[node distance=0pt, inner sep=0pt]
    \node (1) {\strut PRED(};
    \node (2) [right=of 1.east, anchor=west] {\strut $b$};
    \node (3) [right=of 2.east, anchor=west] {\strut )};
    \node (4) [below=4pt of 2.south, anchor=north] {\strut PRED};
    \draw[semithick] ([yshift=2pt]2.south) -- (4.north);
    \node (5) [right=of 4.east, anchor=west] {\strut $^{b}$(${\mathbfit{l}}$)};
\end{tikzpicture} 
\\
\end{tblr} \\

\begin{tikzpicture}
\node[draw=black, minimum height=20pt, minimum width=6.2in] (box) {
{\itshape The book that Lisa {\bfseries bought}\dots}};
\end{tikzpicture} \\*

\begin{tblr}{colspec={X[.6in]X[l]}} 
\textbf{Step 4:} & \textbf{Predicate Algorithm + Argument Addition} \vspace{2pt} \\ 
& Integrate the content of \textit{bought}, represented here as BUY($\alpha$, $\beta$), at the predicate position, extending the argument grid to accommodate a second argument. \vspace{2pt} \\
& 
\begin{tikzpicture}[node distance=0pt, inner sep=0pt]
    \node (1) {\strut PRED(};
    \node (2) [right=of 1.east, anchor=west] {\strut $b$};
    \node (3) [right=of 2.east, anchor=west] {\strut )};
    \node (4) [below=4pt of 2.south, anchor=north] {\bfseries\strut BUY};
    \draw[semithick] ([yshift=2pt]2.south) -- (4.north);
    \node (5) [right=of 4.east, anchor=west] {\strut $^{b}$($l$, ${\mathbfit{\beta}}$)};
\end{tikzpicture} 
\\
\end{tblr} \\

\begin{tblr}{colspec={X[.6in]X[l]}} 
\textbf{Step 5:} & \textbf{Filler Integration + Second-Argument Algorithm} \vspace{2pt} \\ 
& Integrate the filler at the second-argument position. \vspace{2pt} \\
& 
\begin{tikzpicture}[node distance=0pt, inner sep=0pt]
    \node (1) {\strut PRED(};
    \node (2) [right=of 1.east, anchor=west] {\strut $b$};
    \node (3) [right=of 2.east, anchor=west] {\strut )};
    \node (4) [below=4pt of 2.south, anchor=north] {\strut BUY};
    \draw[semithick] ([yshift=2pt]2.south) -- (4.north);
    \node (5) [right=of 4.east, anchor=west] {\strut ($l$, ${\mathbfit{b}}$)};
\end{tikzpicture} 
\\
\end{tblr} \\

\hline
\end{longtblr}
\vspace{-24pt}
\setstretch{1.5}

% Korean

\setstretch{1}
\begin{longtblr}[caption={Processing routine for comprehending sentence-initial \textit{Lisa-{\sc nom} buy-{\sc adn} book\dots} in Korean}, label={tblr:dis-routine-kor}, remark{Note}={PRED~=~predicate position; $\alpha$~=~first-argument position; $\beta$~=~second-argument position; vertical lines indicate modification; superscripts are fillers.}]{colspec={|X[c]|}, rowhead = 0}\hline

\begin{tblr}{colspec={X[.6in]X[l]}} 
\textbf{Step 0:} & \textbf{Semantic Base Projection} \vspace{2pt} \\ 
& Project a minimal predicate-argument complex. \vspace{2pt} \\
& PRED($\alpha$) \\
\end{tblr} \\*

\begin{tikzpicture}
\node[draw=black, minimum height=20pt, minimum width=6.2in] (box) {
{\itshape{\bfseries Lisa-{\sc nom}}\dots\phantom{buy-{\sc adn} book }}};
\end{tikzpicture} \\*

\begin{tblr}{colspec={X[.6in]X[l]}} 
\textbf{Step 1:} & \textbf{First-Argument Algorithm} \vspace{2pt} \\ 
& Integrate the content of \textit{Lisa-{\sc nom}}, represented here as $l$, at the first-argument position.  \vspace{2pt} \\
& PRED($\mathbfit{l}$) \\
\end{tblr} \\*

\begin{tikzpicture}
\node[draw=black, minimum height=20pt, minimum width=6.2in] (box) {
{\itshape Lisa-{\sc nom} {\bfseries buy-{\sc adn}}\dots\phantom{book }}};
\end{tikzpicture} \\*

\begin{tblr}{colspec={X[.6in]X[l]}} 
\textbf{Step 2:} & \textbf{Predicate Algorithm + Argument Addition} \vspace{2pt} \\ 
& Integrate the content of \textit{buy}, represented here as BUY($\alpha$, $\beta$), at the predicate position, expanding the argument grid to accommodate a second argument. \vspace{2pt} \\
& {\bfseries BUY}($l$, $\mathbfit{\beta}$) \\
\end{tblr} \\

\begin{tblr}{colspec={X[.6in]X[l]}} 
\textbf{Step 3:} & \textbf{Relativization Algorithm for Head-Final RCs} \vspace{2pt} \\ 
& Respond to the adnominal marking on \textit{buy} by identifying the missing theme argument as a gap (represented here as $y$) and projecting a new semantic base into which the corresponding filler (also represented as $y$) needs to be integrated. \vspace{2pt} \\
& 
\begin{tikzpicture}[node distance=0pt, inner sep=0pt]
    \node (1) {\bfseries\strut PRED};
    \node (2) [right=of 1.east, anchor=west] {\strut $\mathbfit{^y}$};
    \node (3) [right=of 2.east, anchor=west] {\bfseries\strut ($\mathbfit{\alpha}$)};
    \node (4) [below=4pt of 2.south, anchor=north] {\strut BUY};
    \draw[semithick] ([yshift=5pt]2.south) -- (4.north);
    \node (5) [right=of 4.east, anchor=west] {\strut ($l$, $\mathbfit{y}$)};
\end{tikzpicture} 
\\
\end{tblr} \\

\begin{tikzpicture}
\node[draw=black, minimum height=20pt, minimum width=6.2in] (box) {
{\itshape Lisa-{\sc nom} buy-{\sc adn} {\bfseries book}\dots}};
\end{tikzpicture} \\*

\begin{tblr}{colspec={X[.6in]X[l]}} 
\textbf{Step 4:} & \textbf{Filler Identification} \vspace{2pt} \\ 
& Substitute the content of \textit{book}, represented here as $b$, for $y$. \vspace{2pt} \\
& 
\begin{tikzpicture}[node distance=0pt, inner sep=0pt]
    \node (1) {\strut PRED};
    \node (2) [right=of 1.east, anchor=west] {\strut $\mathbfit{^b}$};
    \node (3) [right=of 2.east, anchor=west] {\strut ($\alpha$)};
    \node (4) [below=4pt of 2.south, anchor=north] {\strut BUY};
    \draw[semithick] ([yshift=6pt]2.south) -- (4.north);
    \node (5) [right=of 4.east, anchor=west] {\strut ($l$, $\mathbfit{b}$)};
\end{tikzpicture} 
\\
\end{tblr} \\

\hline
\end{longtblr}
\vspace{-24pt}
\setstretch{1.5}

% Mandarin

\setstretch{1}
\begin{longtblr}[caption={Processing routine for comprehending sentence-initial \textit{Lisa buy {\sc adn} book\dots} in Mandarin}, label={tblr:dis-routine-zho}, remark{Note}={PRED~=~predicate position; $\alpha$~=~first-argument position; $\beta$~=~second-argument position; vertical lines indicate modification; superscripts are fillers.}]{colspec={|X[c]|}, rowhead = 0}\hline

\begin{tblr}{colspec={X[.6in]X[l]}} 
\textbf{Step 0:} & \textbf{Semantic Base Projection} \vspace{2pt} \\
& Project a minimal predicate-argument complex. \vspace{2pt} \\
& PRED($\alpha$) \\
\end{tblr} \\*

\begin{tikzpicture}
\node[draw=black, minimum height=20pt, minimum width=6.2in] (box) {
{\itshape{\bfseries Lisa}\dots\phantom{buy {\sc adn} book }}};
\end{tikzpicture} \\*

\begin{tblr}{colspec={X[.6in]X[l]}} 
\textbf{Step 1:} & \textbf{First-Argument Algorithm} \vspace{2pt} \\ 
& Integrate the content of \textit{Lisa}, represented here as $l$, at the first-argument position.  \vspace{2pt} \\
& PRED($\mathbfit{l}$) \\
\end{tblr} \\*

\begin{tikzpicture}
\node[draw=black, minimum height=20pt, minimum width=6.2in] (box) {
{\itshape Lisa {\bfseries buy}\dots\phantom{{\sc adn} book }}};
\end{tikzpicture} \\*

\begin{tblr}{colspec={X[.6in]X[l]}} 
\textbf{Step 2:} & \textbf{Predicate Algorithm + Argument Addition} \vspace{2pt} \\ 
& Integrate the content of \textit{buy}, represented here as BUY($\alpha$, $\beta$), at the predicate position, expanding the argument grid to accommodate a second argument. \vspace{2pt} \\
& {\bfseries BUY}($l$, $\mathbfit{\beta}$) \\
\end{tblr} \\

\begin{tikzpicture}
\node[draw=black, minimum height=20pt, minimum width=6.2in] (box) {
{\itshape Lisa buy {\bfseries{\sc adn}}\dots\phantom{book }}};
\end{tikzpicture} \\*

\begin{tblr}{colspec={X[.6in]X[l]}} 
\textbf{Step 3:} & \textbf{Relativization Algorithm for Head-Final RCs} \vspace{2pt} \\ 
& Respond to the adnominal particle by identifying the missing theme argument as a gap (represented here as $y$) and projecting a new semantic base into which the corresponding filler (also represented as $y$) needs to be integrated. \vspace{2pt} \\
& 
\begin{tikzpicture}[node distance=0pt, inner sep=0pt]
    \node (1) {\bfseries\strut PRED};
    \node (2) [right=of 1.east, anchor=west] {\strut $\mathbfit{^y}$};
    \node (3) [right=of 2.east, anchor=west] {\bfseries\strut ($\mathbfit{\alpha}$)};
    \node (4) [below=4pt of 2.south, anchor=north] {\strut BUY};
    \draw[semithick] ([yshift=5pt]2.south) -- (4.north);
    \node (5) [right=of 4.east, anchor=west] {\strut ($l$, $\mathbfit{y}$)};
\end{tikzpicture} 
\\
\end{tblr} \\

\begin{tikzpicture}
\node[draw=black, minimum height=20pt, minimum width=6.2in] (box) {
{\itshape Lisa buy {\sc adn} {\bfseries book}\dots}};
\end{tikzpicture} \\*

\begin{tblr}{colspec={X[.6in]X[l]}} 
\textbf{Step 4:} & \textbf{Filler Identification} \vspace{2pt} \\ 
& Substitute the content of \textit{book}, represented here as $b$, for $y$. \vspace{2pt} \\
& 
\begin{tikzpicture}[node distance=0pt, inner sep=0pt]
    \node (1) {\strut PRED};
    \node (2) [right=of 1.east, anchor=west] {\strut $\mathbfit{^b}$};
    \node (3) [right=of 2.east, anchor=west] {\strut ($\alpha$)};
    \node (4) [below=4pt of 2.south, anchor=north] {\strut BUY};
    \draw[semithick] ([yshift=6pt]2.south) -- (4.north);
    \node (5) [right=of 4.east, anchor=west] {\strut ($l$, $\mathbfit{b}$)};
\end{tikzpicture} 
\\
\end{tblr} \\

\hline
\end{longtblr}
\vspace{-6pt}
\setstretch{1.5}

To make the routines as simple as possible, I placed the RC and its head noun in sentence-initial position in Tables~\ref{tblr:dis-routine-eng}--\ref{tblr:dis-routine-zho}; things would play out slightly differently if I put them elsewhere in the sentence. In Korean, the head noun can carry case marking which would allow the filler to be integrated instantly into the new predicate-argument complex; however, I decided to show the head noun without case marking here to keep the routines for Korean and Mandarin maximally similar.

For comparison, let us also consider the case of Hebrew, which like English has filler-gap RCs. The Hebrew translation equivalent of \textit{the~book that Lisa bought} is given in (\ref{ex-dis-heb}), and a processing routine for comprehending it is presented in Table~\ref{tblr:dis-routine-heb}. Notice that the sequence of operations in the routine for this Hebrew RC is identical to that for its English translation equivalent in Table~\ref{tblr:dis-routine-eng}.

% Hebrew

\ex
\begingl 
\gla ha-sefer she-Lisa kanta // \label{ex-dis-heb}
\glb the-book that-Lisa bought //
\glft `the book that Lisa bought' //
\endgl\par\nobreak\vspace{5pt}
\hfill (based on \cite{koesterich.meltzer-asscher2022}, Slide 3, (1))
\vspace{-12pt}
\xe 

\setstretch{1}
\begin{longtblr}[caption={Processing routine for comprehending sentence-initial \textit{The-book that-Lisa bought\dots} in Hebrew}, label={tblr:dis-routine-heb}, remark{Note}={PRED~=~predicate position; $\alpha$~=~first-argument position; $\beta$~=~second-argument position; vertical lines indicate modification; superscripts are fillers.}]{colspec={|X[c]|}, rowhead = 0}\hline

\begin{tblr}{colspec={X[.6in]X[l]}} 
\textbf{Step 0:} & \textbf{Semantic Base Projection} \vspace{2pt} \\ 
& Project a minimal predicate-argument complex. \vspace{2pt} \\
& PRED($\alpha$) \\
\end{tblr} \\*

\begin{tikzpicture}
\node[draw=black, minimum height=20pt, minimum width=6.2in] (box) {
{\itshape{\bfseries The-book}\dots\phantom{that-Lisa bought }}};
\end{tikzpicture} \\*

\begin{tblr}{colspec={X[.6in]X[l]}} 
\textbf{Step 1:} & \textbf{First-Argument Algorithm} \vspace{2pt} \\ 
& Integrate the content of \textit{the-book}, represented here as $b$, at the first-argument position. \vspace{2pt} \\
& PRED($\mathbfit{b}$) \\
\end{tblr} \\

\begin{tikzpicture}
\node[draw=black, minimum height=20pt, minimum width=6.2in] (box) {
{\itshape The-book {\bfseries that-Lisa}\dots\phantom{bought }}};
\end{tikzpicture} \\*

\begin{tblr}{colspec={X[.6in]X[l]}} 
\textbf{Step 2:} & \textbf{Relativization Algorithm for Head-Initial RCs} \vspace{2pt} \\ 
& Respond to the relativizer \textit{she-} by projecting a semantic base as a modifier of $b$ and give it a filler, represented here as a superscript $b$, that needs to be integrated into its argument grid. \vspace{2pt} \\
& 
\begin{tikzpicture}[node distance=0pt, inner sep=0pt]
    \node (1) {\strut PRED(};
    \node (2) [right=of 1.east, anchor=west] {\strut $b$};
    \node (3) [right=of 2.east, anchor=west] {\strut )};
    \node (4) [below=4pt of 2.south, anchor=north] {\bfseries\strut PRED};
    \draw[semithick] ([yshift=2pt]2.south) -- (4.north);
    \node (5) [right=of 4.east, anchor=west] {\bfseries\strut $^{\mathbfit{b}}$(${\mathbfit{\alpha}}$)};
\end{tikzpicture} 
\\
\end{tblr} \\

\begin{tblr}{colspec={X[.6in]X[l]}} 
\textbf{Step 3:} & \textbf{First-Argument Algorithm} \vspace{2pt} \\ 
& Integrate the content of \textit{Lisa}, represented here as $l$, at the first-argument position. \vspace{2pt} \\
& 
\begin{tikzpicture}[node distance=0pt, inner sep=0pt]
    \node (1) {\strut PRED(};
    \node (2) [right=of 1.east, anchor=west] {\strut $b$};
    \node (3) [right=of 2.east, anchor=west] {\strut )};
    \node (4) [below=4pt of 2.south, anchor=north] {\strut PRED};
    \draw[semithick] ([yshift=2pt]2.south) -- (4.north);
    \node (5) [right=of 4.east, anchor=west] {\strut $^{b}$(${\mathbfit{l}}$)};
\end{tikzpicture} 
\\
\end{tblr} \\

\begin{tikzpicture}
\node[draw=black, minimum height=20pt, minimum width=6.2in] (box) {
{\itshape The-book that-Lisa {\bfseries bought}\dots}};
\end{tikzpicture} \\*

\begin{tblr}{colspec={X[.6in]X[l]}} 
\textbf{Step 4:} & \textbf{Predicate Algorithm + Argument Addition} \vspace{2pt} \\ 
& Integrate the content of \textit{bought}, represented here as BUY($\alpha$, $\beta$), at the predicate position, extending the argument grid to accommodate a second argument. \vspace{2pt} \\
& 
\begin{tikzpicture}[node distance=0pt, inner sep=0pt]
    \node (1) {\strut PRED(};
    \node (2) [right=of 1.east, anchor=west] {\strut $b$};
    \node (3) [right=of 2.east, anchor=west] {\strut )};
    \node (4) [below=4pt of 2.south, anchor=north] {\bfseries\strut BUY};
    \draw[semithick] ([yshift=2pt]2.south) -- (4.north);
    \node (5) [right=of 4.east, anchor=west] {\strut $^{b}$($l$, ${\mathbfit{\beta}}$)};
\end{tikzpicture} 
\\
\end{tblr} \\

\begin{tblr}{colspec={X[.6in]X[l]}} 
\textbf{Step 5:} & \textbf{Filler Integration + Second-Argument Algorithm} \vspace{2pt} \\ 
& Integrate the filler at the second-argument position. \vspace{2pt} \\
& 
\begin{tikzpicture}[node distance=0pt, inner sep=0pt]
    \node (1) {\strut PRED(};
    \node (2) [right=of 1.east, anchor=west] {\strut $b$};
    \node (3) [right=of 2.east, anchor=west] {\strut )};
    \node (4) [below=4pt of 2.south, anchor=north] {\strut BUY};
    \draw[semithick] ([yshift=2pt]2.south) -- (4.north);
    \node (5) [right=of 4.east, anchor=west] {\strut ($l$, ${\mathbfit{b}}$)};
\end{tikzpicture} 
\\
\end{tblr} \\

\hline
\end{longtblr}
\vspace{-6pt}
\setstretch{1.5}

As a follow-up to the present dissertation research, future studies could investigate the degree to which these sorts of comparisons affect the L2 processing of RCs. A close to ideal pairing would be L1-Mandarin L2ers of English (MLEs) vs. L1-Hebrew L2ers of English (HLEs). Mandarin and Hebrew share the same distribution of grammatical gapping and resumption (for Mandarin, see \cite{keenan.comrie1977}, p.~93; \cite{li.wei2014}, p.~286; \cite{pan2016}, p.~292ff; for Hebrew, see \cite{borer1984}, p.~220; \cite{fadlon.etal2019}, p.~46; \cite{keenan.comrie1977}, p.~93), allowing or requiring RPs in many RC environments where English prohibits them. As we have seen in this section, however, these two languages differ with respect to how similar their RC processing routines are to those used in English. Given the high level of overlap between the processing routines for many RC types in English and Hebrew, we might expect that in the L2 processing of English RCs, HLEs would experience less difficulty compared to MLEs, especially at lower levels of English proficiency. It is also conceivable that the degree of overlap between routines for processing RCs in the L1 (viz.,~Hebrew or Mandarin) vs. the target language (viz.,~English) might affect the rates at which L2ers resort to resumptive RCs while attempting to produce English sentences.\footnotemark{}

\footnotetext{In an investigation such as the one I am proposing here, it would be important to take into consideration the relative frequency of resumptive RCs in spoken/written Hebrew vs. Mandarin in each of the relevant syntactic environments, since asymmetries in this regard could influence participants' performance in the tasks.}

If it is not possible to find HLEs with sufficiently low English proficiency and/or age of onset for English that is sufficiently late, then L1-Arabic L2ers could potentially be tested instead. The distribution of grammatical resumption in Arabic closely resembles that in both Hebrew and Mandarin (\cite{algady2013}, p.~45; \cite{keenan.comrie1977}, p.~93; \cite{maghrabi1997}, p.~89), and Arabic is also similar to English and Hebrew in having postnominal RCs and optionality between SVO and VSO word order (Maghrabi, 1997, p. 81ff). 

%========================================================%
\section{Summing up}\label{dis-conclusion}
%========================================================%

This dissertation examined the processing and acceptability of gapped vs. resumptive ORCs/SRCs among L1-English controls, L1-Korean L2ers of English, and L1-Mandarin L2ers of English (all of them adults). The main tasks consisted of an oral elicited production task probing processing during RC production, a self-paced reading task probing processing during RC comprehension, and an acceptability judgment task (with versions in English, Korean, and Mandarin) testing the offline acceptability of the sentence types used in the other two tasks. The conditions in each of these tasks targeted three increasingly complex RC environments: short-distance, long-distance, and \textit{wh}-island. The results provide clear-cut evidence that resumption facilitates both production and comprehension of RCs under processing strain---for native speakers as well as for L2ers. A sizable portion of the participants in each group showed signs of deriving a processing benefit from resumption even when they did not consistently accept RPs in the offline judgment task, suggesting that English resumptive RCs for them are processing resumptives rather than grammatical resumptives. Furthermore, higher L2 proficiency in English was associated with higher ratings for gaps and lower ratings for RPs, at least in domains where gapped RCs are permitted in English, indicating that L2ers are able to develop a sensitivity to the English prohibition against resumption in RCs; this sensitivity, moreover, appears to develop faster where processing is easier. 

There are a number of different matters not directly addressed in this dissertation that I see as priorities for future research. First, it will be important to investigate how L2ers ostensibly come to have resumptive RCs that arise out of processing. Second, future studies should examine other language pairings in order to examine how differences---specifically, crosslinguistic differences and processing-routine differences---between the L1 vs. the target language affect the patterns that L2ers evince when producing and comprehending RCs. Third, it will be necessary to explore whether findings of the type obtained in this study extend to other types of syntactic phenomena that emerge in interlanguages despite being unacceptable in both the L1 and the target languages: Examples include \textit{wh}-scope marking in long-distance questions (as in \textit{Who does Mary believe who John saw yesterday}; e.g.,~\cite{schulz2006}) and subject-auxiliary inversion in embedded questions (as in \textit{Mary remembered what will John do tomorrow}; e.g.,~\cite{pozzan.quirk2014}). Finally, the approach adopted in this project could be leveraged to probe other processing hierarchies, such as Rizzi's (\citeyear{rizzi2018}) distinctness hierarchy for intervention effects. Continued research in these areas has the potential to advance our understanding of the interplay between language competence and language processing during (L1 and) L2 acquisition.